\documentclass[aspectratio=169,11pt]{beamer}

% Theme and color setup
\usetheme{default}
\usecolortheme{default}

% Remove navigation symbols
\setbeamertemplate{navigation symbols}{}

% Custom colors (NYU-inspired: purple accent)
\definecolor{nyupurple}{RGB}{87,35,100}
\definecolor{darkblue}{RGB}{0,40,85}
\definecolor{lightgray}{RGB}{245,245,245}
\definecolor{darkgray}{RGB}{64,64,64}

% Set main colors
\setbeamercolor{frametitle}{fg=white,bg=darkblue}
\setbeamercolor{title}{fg=darkblue}
\setbeamercolor{author}{fg=black}
\setbeamercolor{institute}{fg=black}
\setbeamercolor{date}{fg=black}
\setbeamercolor{structure}{fg=darkblue}
\setbeamercolor{normal text}{fg=black}
\setbeamercolor{itemize item}{fg=nyupurple}
\setbeamercolor{itemize subitem}{fg=darkblue}
\setbeamercolor{enumerate item}{fg=darkblue}

% Frame title band
\newlength{\titlebandht}
\setlength{\titlebandht}{4ex}

\setbeamertemplate{frametitle}{
  \nointerlineskip
  \begin{beamercolorbox}[wd=\paperwidth,ht=\titlebandht,dp=1.1ex,leftskip=1.2em,rightskip=1.2em]{frametitle}
    \usebeamerfont{frametitle}\insertframetitle
    \ifx\insertframesubtitle\@empty\relax\else\\[-0.2ex]
      \usebeamerfont{framesubtitle}\insertframesubtitle
    \fi
  \end{beamercolorbox}
  \vspace{0.6em}
}

% Font settings
\usefonttheme{professionalfonts}
\usefonttheme{serif}

% Footline with page numbers
\setbeamertemplate{footline}{
    \hfill\insertframenumber\hspace{2em}\vspace{0.5em}
}

% Packages
\usepackage{amsmath,amssymb,amsthm}
\usepackage{graphicx}
\usepackage{tikz}
\usepackage{booktabs}
\usepackage{mathtools}
\usepackage{bm}
\usepackage{bbm}

% Custom commands for math
\newcommand{\E}{\mathbb{E}}
\renewcommand{\P}{\mathbb{P}}
\newcommand{\R}{\mathbb{R}}
\newcommand{\N}{\mathbb{N}}
\newcommand{\Z}{\mathbb{Z}}
\newcommand{\F}{\mathcal{F}}
\newcommand{\Cov}{\text{Cov}}
\newcommand{\Var}{\text{Var}}
\newcommand{\Corr}{\text{Corr}}
\newcommand{\xbf}{\mathbf{x}}
\newcommand{\ybf}{\mathbf{y}}
\newcommand{\zbf}{\mathbf{z}}
\newcommand{\ubf}{\mathbf{u}}
\newcommand{\vbf}{\mathbf{v}}
\newcommand{\hbf}{\mathbf{h}}
\DeclareMathOperator*{\argmin}{arg\,min}
\DeclareMathOperator*{\argmax}{arg\,max}
\DeclareMathOperator*{\plim}{plim}

% Block colors
\setbeamercolor{block title}{bg=lightgray,fg=darkblue}
\setbeamercolor{block body}{bg=lightgray!50}
\setbeamercolor{block title alerted}{bg=nyupurple!15,fg=nyupurple}
\setbeamercolor{block body alerted}{bg=nyupurple!5}

% Figure path
\graphicspath{{../tex/}}

% Title information
\title{Problem Set 2: Interpolation, Occupational Choice, and Search Frictions}
\subtitle{Solutions + Comments}
\author{Rafael Lincoln}
\institute{New York University}
\date{\today}

\begin{document}

% ============================================================================
% SLIDE 1: TITLE
% ============================================================================
\begin{frame}[plain]
    \titlepage
\end{frame}

% ============================================================================
%  PROBLEM 1
% ============================================================================

% ============================================================================
% SLIDE 2: 1D INTERPOLATION -- SETUP
% ============================================================================
\begin{frame}{Problem 1: 1D Interpolation --- Setup}
    \textbf{Goal}: Approximate three benchmark functions on $s\in[0,1]$ and compare methods.
    \vspace{0.3em}

    \textbf{Functions} (chosen for difficulty):
    \begin{itemize}\setlength\itemsep{0pt}
        \item $f_1(s) = \exp(-s/0.01)$ \hfill \textit{steep near $s=0$}
        \item $f_2(s) = \max(s, 0.5)$ \hfill \textit{kink at $s=0.5$}
        \item $f_3(s) = \sin(25s^2) + \tfrac{1}{2}\log(s+0.01)$ \hfill \textit{oscillatory + log}
    \end{itemize}
    \vspace{0.3em}

    \textbf{Methods}: 
    \begin{itemize}\setlength\itemsep{0pt}
        \item \textbf{Chebyshev} polynomial via least-squares fit
        \item \textbf{Linear spline} ($k=1$, uniform grid)
        \item \textbf{Cubic spline} ($k=3$, uniform grid)
    \end{itemize}
    
    \textbf{Error metrics} (on 10\,001-point fine grid):
    \[
    \text{Max error} = \max_{s\in\mathcal{S}_{\text{fine}}} |\hat f(s) - f(s)|,
    \qquad
    \text{Mean error} = \frac{1}{|\mathcal{S}_{\text{fine}}|}\sum_s |\hat f(s) - f(s)|
    \]
\end{frame}

% ============================================================================
% SLIDE 3: 1D INTERPOLATION -- RESULTS
% ============================================================================
\begin{frame}{Problem 1: 1D Interpolation --- Results}
    \vspace{-0.6em}
    \begin{columns}[T]
        \begin{column}{0.50\textwidth}
            \begin{center}
                \includegraphics[height=0.66\textheight]{fig_problem1_1d_interpolation.pdf}
            \end{center}
        \end{column}
        \begin{column}{0.48\textwidth}
            \footnotesize
            \textbf{Error summary ($n=51$)}\\[-0.2em]
            \textit{max / mean}\\[0.4em]
            \begin{itemize}\setlength\itemsep{0.35em}
                \item \textbf{$f_1$ (exp)}\\
                Cheb: \texttt{6.5e-12 / 1.6e-12}\\
                Linear: \texttt{2.1e-01 / 3.1e-03}\\
                Cubic: \texttt{1.4e-01 / 2.4e-03}
                \item \textbf{$f_2$ (kink)}\\
                Cheb: \texttt{2.9e-03 / 3.7e-04}\\
                Linear: \texttt{1.1e-16 / 1.5e-17}\\
                Cubic: \texttt{1.7e-03 / 5.8e-05}
                \item \textbf{$f_3$ (sin+log)}\\
                Cheb: \texttt{6.8e-06 / 2.3e-07}\\
                Linear: \texttt{1.2e-01 / 1.9e-02}\\
                Cubic: \texttt{5.0e-02 / 1.1e-03}
            \end{itemize}
        \end{column}
    \end{columns}
    \vspace{0.1em}
    \footnotesize
    \textbf{Takeaway}: Chebyshev excels on smooth targets; cubic spline handles kinks better. Linear spline converges slowly but never overshoots.
\end{frame}

% ============================================================================
% SLIDE 4: OCCUPATIONAL CHOICE -- ECONOMIC SETTING
% ============================================================================
\begin{frame}{Problem 1.2: Occupational Choice --- Setting}
    \small
    \textbf{Agents}: Infinitely-lived households with assets $a\ge 0$, entrepreneurial ability $z$, labor ability $e$.

    \textbf{Timing}: Observe $(a_t,e_t,z_t)$ $\to$ choose occupation $\to$ consume $c_t$, save $a_{t+1}\ge 0$.
    \vspace{0.3em}

    \begin{columns}[T]
        \begin{column}{0.45\textwidth}
            \textbf{Worker}:
            \begin{itemize}\setlength\itemsep{0pt}
                \item Income $y^W = We$
                \item No capital needed
            \end{itemize}
        \end{column}
        \begin{column}{0.50\textwidth}
            \textbf{Entrepreneur}:
            \begin{itemize}\setlength\itemsep{0pt}
                \item Technology $z(k^\alpha \ell^{1-\alpha})^\eta$
                \item Collateral constraint: $k \le \lambda a$
                \item Hires labor at $W$, rents capital at $R$
            \end{itemize}
        \end{column}
    \end{columns}
    \vspace{0.3em}
    \begin{alertblock}{Key friction}
        \small The collateral constraint $k\le\lambda a$ prevents small-wealth, high-ability agents from becoming entrepreneurs even when it would be efficient.
    \end{alertblock}
\end{frame}

% ============================================================================
% SLIDE 5: OCCUPATIONAL CHOICE -- KEY EQUATIONS
% ============================================================================
\begin{frame}{Problem 1.2: Key Equations}
    \textbf{Entrepreneur profit} (after substituting optimal labor):
    \[
    \pi(a,z) = \max_{k \le \lambda a}\; \tilde\pi(k,z),
    \qquad
    k(a,z) = \min\{k^u(z),\; \lambda a\}
    \]
    where $k^u(z)$ is the unconstrained optimum.
    \vspace{0.4em}

    \textbf{Occupation value}:
    \[
    V(a,e,z) = \max\bigl\{\pi(a,z),\; We\bigr\}
    \]
    \vspace{0.2em}

    \textbf{Bellman equation} (dynamic extension):
    \[
    \boxed{
    J(a,e,z) = \max_{a'\ge 0}\;\bigl\{u\bigl((1+r)a + V(a,e,z) - a'\bigr) + \beta\,\E\bigl[J(a',e',z')\mid e,z\bigr]\bigr\}
    }
    \]
    \vspace{0.2em}

    \textbf{Problem 1.2 specifically}: approximate $V(a,e,z)$ on a 3D grid using Chebyshev, linear spline, and cubic spline --- same comparison as 1D but in three dimensions.
\end{frame}

% ============================================================================
% SLIDE 6: 3D APPROXIMATION OF V
% ============================================================================
\begin{frame}{Problem 1.2: 3D Approximation of $V(a,e,z)$}
    \vspace{-0.6em}
    \begin{columns}[T]
        \begin{column}{0.52\textwidth}
            \footnotesize
            \textbf{Setup}: tensor grid over $(a,z,e)$ with $n$ points per dimension; errors on a $(3n)^3$ fine grid.
            \vspace{0.3em}

            \textbf{Numerical story}:
            \begin{itemize}\setlength\itemsep{0.25em}
                \item \textbf{Chebyshev basis + imposed regular grid} $\Rightarrow$ ill-conditioned \textbf{Vandermonde} design matrix; coefficients become unstable as $n$ rises.
                \item \textbf{Splines} are local and designed for regular grids $\Rightarrow$ stable, smooth error decay.
                \item Using \textbf{Chebyshev nodes} (instead of imposed nodes) would greatly improve conditioning.
            \end{itemize}
        \end{column}
        \begin{column}{0.46\textwidth}
            \scriptsize
            \textbf{Key numbers (mean abs.\ error)}\\[-0.2em]
            \textit{($n$ points per dimension)}\\[0.3em]
            \begin{center}
            \renewcommand{\arraystretch}{1.15}
            \begin{tabular}{lccc}
                \toprule
                & $n=11$ & $n=15$ & $n=21$ \\
                \midrule
                Cheb & 0.135 & 0.201 & 53.1 \\
                Linear & 0.210 & 0.112 & 0.057 \\
                Cubic & 0.083 & 0.045 & 0.022 \\
                \bottomrule
            \end{tabular}
            \end{center}
            \vspace{0.2em}
            \scriptsize
            \textbf{Cheb max abs.\ error explodes}: $8.5 \to 166 \to 6.4\times 10^{5}$.
        \end{column}
    \end{columns}
    \vspace{0.1em}
    \footnotesize
    \textbf{Takeaway}: On the imposed grid, splines convey the approximation signal cleanly; Chebyshev is dominated by conditioning unless we use Chebyshev nodes.
\end{frame}

% ============================================================================
% SLIDE 7: DYNAMIC EXTENSION -- VFI + STATIONARY DISTRIBUTION
% ============================================================================
\begin{frame}{Dynamic Extension: VFI + Stationary Distribution}
    \textbf{Calibration highlights}:
    \begin{itemize}
        \item $\beta = 0.95$, $\gamma = 2$ (CRRA), $r = 0.02$
        \item Abilities: $\log e$, $\log z$ each AR(1), Rouwenhorst ($n_e=n_z=5$; $\rho=0.9$, $\sigma_e=0.2$, $\sigma_z=0.3$)
        \item Joint transition: $\Pi = \Pi_e \otimes \Pi_z$
    \end{itemize}
    \vspace{0.3em}

    \textbf{Value Function Iteration}:
    \begin{enumerate}
        \item For each state $(a,e,z)$: compute cash-on-hand $(1+r)a + V(a,e,z)$
        \item Maximize over $a'\ge 0$ on a dense choice grid; evaluate continuation values by \textbf{interpolating on the $a$-grid} (1D interpolation in $a'$ of $\beta\,\E[J(a',e',z')]$)
    \end{enumerate}
    \vspace{0.3em}

    \textbf{Stationary distribution} via forward iteration:
    \begin{enumerate}
        \item \textbf{Policy step}: move mass according to $a'(a,e,z)$, using lottery weights for off-grid $a'$
        \item \textbf{Exogenous step}: apply $(\Pi_e\otimes\Pi_z)^\top$
        \item Iterate until $\|\mu^{(k+1)} - \mu^{(k)}\| < \varepsilon$
    \end{enumerate}
\end{frame}

% ============================================================================
% SLIDE 8: EXTENSION RESULTS -- POLICIES
% ============================================================================
\begin{frame}{Extension Results: Policy Functions}
    \vspace{-0.3em}
    \begin{center}
        \includegraphics[height=0.72\textheight]{fig_problem1_extension_policies.pdf}
    \end{center}
    \vspace{-0.8em}
    \small
    \textbf{Interpretation}: Savings $a'-a$ increasing in $a$; entrepreneurs save more. Consumption smooth and concave.
\end{frame}

% ============================================================================
% SLIDE 9: EXTENSION RESULTS -- OCCUPATIONAL CHOICE MAP
% ============================================================================
\begin{frame}{Extension Results: Occupational Choice}
    \begin{center}
        \includegraphics[width=0.88\textwidth]{fig_problem1_extension_occ_map.pdf}
    \end{center}
    \vspace{-0.5em}
    \small
    \textbf{Interpretation}: Wealthier and higher-$z$ agents become entrepreneurs (collateral constraint relaxes). The cutoff $a^*(z;e)$ shifts left as $z$ rises --- high-ability agents need less wealth to find entrepreneurship profitable.
\end{frame}

% ============================================================================
% SLIDE 10: EXTENSION RESULTS -- WEALTH DISTRIBUTION
% ============================================================================
\begin{frame}{Extension Results: Wealth Concentration}
    \begin{center}
        \includegraphics[width=0.88\textwidth]{fig_problem1_extension_lorenz_ccdf.pdf}
    \end{center}
    \vspace{-0.5em}
    \small
    \textbf{Interpretation}: The Lorenz curve shows significant inequality; the CCDF exhibits a Pareto-like tail. Entrepreneurial returns --- amplified by the collateral constraint --- drive right-tail wealth concentration: those who can invest earn disproportionately high returns, accumulate faster, and stay rich.
\end{frame}

% ============================================================================
%  PROBLEM 2
% ============================================================================

% ============================================================================
% SLIDE 11: DMP -- ENVIRONMENT
% ============================================================================
\begin{frame}{Problem 2: DMP Model --- Environment}
    \footnotesize
    \vspace{-0.5em}
    \begin{columns}[T]
        \begin{column}{0.56\textwidth}
            \textbf{Agents / states}
            \begin{itemize}\setlength\itemsep{0.15em}
                \item \textbf{Workers}: employed or unemployed.
                \item \textbf{Firms}: post vacancies at flow cost $\kappa$; if matched, produce.
                \item \textbf{State}: aggregate productivity $z_t$.
            \end{itemize}
            \vspace{0.15em}

            \textbf{Within-period timing}
            \begin{itemize}\setlength\itemsep{0.15em}
                \item Produce $\to$ separation (rate $\sigma$) $\to$ matching $\to$ Nash bargaining.
                \item Productivity (AR(1) in logs): $\log z_{t+1}=\rho_z\log z_t+\sigma_z\varepsilon_{t+1}$.
            \end{itemize}
        \end{column}
        \begin{column}{0.42\textwidth}
            \textbf{Matching technology}
            \begin{itemize}\setlength\itemsep{0.15em}
                \item Matching: $m_t = B\,u_t^\alpha v_t^{1-\alpha}$.
                \item Tightness: $\theta_t \equiv v_t/u_t$.
                \item Job finding: $\lambda_{w,t}=m_t/u_t=B\theta_t^{1-\alpha}$.
                \item Vacancy filling: $\lambda_{f,t}=m_t/v_t=B\theta_t^{-\alpha}$.
            \end{itemize}
            \vspace{0.15em}

            \textbf{Parameter map}
            \begin{itemize}\setlength\itemsep{0.15em}
                \item $B,\alpha$: matching efficiency / elasticity.
                \item $\sigma,\kappa$: separation / vacancy cost.
                \item $b,\beta,\gamma$: outside value / discount / bargaining.
                \item $\rho_z,\sigma_z$: persistence / volatility.
            \end{itemize}
        \end{column}
    \end{columns}
\end{frame}

% ============================================================================
% SLIDE 12: DMP -- VALUE FUNCTIONS AND SURPLUS
% ============================================================================
\begin{frame}{Problem 2: Value Functions and Surplus}
    \footnotesize
    \vspace{-0.5em}
    \begin{columns}[T]
        \begin{column}{0.53\textwidth}
            \textbf{Objects (what they mean)}
            \begin{itemize}\setlength\itemsep{0.15em}
                \item $W(z)$ / $U(z)$: worker value employed / unemployed (flow $b$).
                \item $J(z)$: firm value from a filled job (match).
                \item $S(z)\equiv W(z)-U(z)+J(z)$: total match surplus.
            \end{itemize}
            \vspace{0.15em}

            \textbf{Bellman equations (compact)}
            \[
            \begin{aligned}
            W(z)&= w(z)+\beta\,\E[(1-\sigma)W(z')+\sigma U(z')],\\
            U(z)&= b+\beta\,\E[\lambda_w(z)W(z')+(1-\lambda_w(z))U(z')],\\
            J(z)&= z-w(z)+\beta(1-\sigma)\E[J(z')].
            \end{aligned}
            \]
        \end{column}
        \begin{column}{0.44\textwidth}
            \textbf{Mapping into equilibrium conditions}
            \begin{itemize}\setlength\itemsep{0.15em}
                \item \textbf{Nash bargaining} (share $\gamma$):\ $W-U=\gamma S$, $J=(1-\gamma)S$.
                \item \textbf{Surplus recursion} (substitute out $w$):
            \end{itemize}
            \[
            \boxed{\begin{aligned}
            S(z) &= z-b \\
            &\quad + \beta\bigl(1-\sigma-\gamma\lambda_w(z)\bigr)\E[S(z')\mid z]
            \end{aligned}}
            \]
            \vspace{0.1em}
            \textbf{Free entry} (vacancy value $=0$) $\Rightarrow$
            \[
            \boxed{\lambda_w(z)=\chi\cdot\bigl(\beta\,\E[S(z')\mid z]\bigr)^{\frac{1-\alpha}{\alpha}}}
            \]
            {\scriptsize where $\chi \equiv B^{1/\alpha}\Bigl(\frac{1-\gamma}{\kappa}\Bigr)^{\frac{1-\alpha}{\alpha}}$.}
        \end{column}
    \end{columns}
\end{frame}

% ============================================================================
% SLIDE 13: DMP -- EQUILIBRIUM AND CALIBRATION
% ============================================================================
\begin{frame}{Problem 2: Equilibrium and Calibration}
    \textbf{Recursive equilibrium}: two equations in $(S, \lambda_w)$ indexed by $z$-states.
    \vspace{0.3em}

    \begin{center}
    \small
    \begin{tabular}{llcl}
        \toprule
        \textbf{Parameter} & \textbf{Symbol} & \textbf{Value} & \textbf{Description} \\
        \midrule
        Discount factor & $\beta$ & $0.96^{1/12}$ & Monthly \\
        Separation & $\sigma$ & $0.04$ & Monthly rate \\
        Bargaining & $\gamma$ & $0.50$ & Worker share \\
        Outside option & $b$ & $0.70$ & Flow value \\
        Persistence & $\rho_z$ & $0.90^{1/12}$ & AR(1) monthly \\
        Volatility & $\sigma_z$ & $0.005$ & Innovation std.\ dev. \\
        Matching elast. & $\alpha$ & $0.50$ & CRS exponent \\
        Matching eff. & $B$ & $0.40$ & Scale \\
        Vacancy cost & $\kappa$ & implied & Targets $\bar\theta=1$ \\
        \bottomrule
    \end{tabular}
    \end{center}
    \vspace{0.3em}
    \small
    Discretize $z$ with Rouwenhorst ($N=51$ states). Steady state: $\bar\lambda_w = B = 0.4$, $\bar\theta = 1$.
\end{frame}

% ============================================================================
% SLIDE 14: DMP -- SOLUTION METHODS
% ============================================================================
\begin{frame}{Problem 2: Solution Methods}
    \textbf{Three approaches} to solve the discretized system:
    \vspace{0.3em}

    \begin{enumerate}
        \item \textbf{Perturbation (Dynare, order=1)}: log-linearize around steady state
        \[
        \tilde s_t = \frac{1}{\bar S}\hat z_t + \beta\Bigl(1-\sigma-\frac{\gamma}{\alpha}\bar\lambda_w\Bigr)\E_t[\tilde s_{t+1}],
        \quad
        \tilde\lambda_t = \frac{1-\alpha}{\alpha}\E_t[\tilde s_{t+1}]
        \]
        \item \textbf{Iterative fixed-point}: given $\lambda_w^{(k)}$, solve linear system for $S^{(k)}$; update $\lambda_w^{(k+1)}$
        \item \textbf{Newton / Broyden}: stack $x = (S, \lambda_w) \in \R^{2N}$; solve $F(x)=0$ with exact Jacobian or quasi-Newton
    \end{enumerate}
    \vspace{0.4em}

    \begin{center}
    \small
    \begin{tabular}{lcc}
        \toprule
        Method & Iterations & Agree? \\
        \midrule
        Iterative & $\sim$125 & \checkmark \\
        Broyden & $\sim$15 & \checkmark \\
        Newton (exact) & $\sim$6 & \checkmark \\
        \bottomrule
    \end{tabular}
    \end{center}
\end{frame}

% ============================================================================
% SLIDE 15: DMP -- POLICY FUNCTIONS
% ============================================================================
\begin{frame}{Problem 2: Policy Functions $S(z)$ and $\lambda_w(z)$}
    \begin{center}
        \includegraphics[width=0.92\textwidth]{fig_problem2_profiles.pdf}
    \end{center}
    \vspace{-0.5em}
    \small
    All three nonlinear methods produce \textbf{identical} policy functions. Both $S(z)$ and $\lambda_w(z)$ are increasing in $z$: higher productivity makes matches more valuable and induces more vacancy posting.
\end{frame}

% ============================================================================
% SLIDE 16: DMP -- DYNARE VS NONLINEAR
% ============================================================================
\begin{frame}{Problem 2: Dynare Linear vs.\ Nonlinear Solution}
    \begin{center}
        \includegraphics[width=0.92\textwidth]{fig_problem2_profiles_dynare_linear_compare.pdf}
    \end{center}
    \vspace{-0.5em}
    \small
    The first-order perturbation (Dynare) is \textbf{accurate near the steady state} ($z=1$) but \textbf{underestimates the response} at tail values of $z$. The nonlinear solution captures the curvature that linearization misses.
\end{frame}

% ============================================================================
% SLIDE 17: DMP -- SIMULATED FLUCTUATIONS
% ============================================================================
\begin{frame}{Problem 2: Simulated Labor-Market Fluctuations}
    \vspace{-0.3em}
    \begin{center}
        \includegraphics[height=0.70\textheight]{fig_problem2_simulated_labor_fluctuations.pdf}
    \end{center}
    \vspace{-0.8em}
    \small
    \textbf{Amplification}: $z\!\uparrow$ $\Rightarrow$ $S\!\uparrow$ $\Rightarrow$ vacancy posting $\uparrow$ $\Rightarrow$ $\theta\!\uparrow$, $\lambda_w\!\uparrow$, $v\!\uparrow$; unemployment falls. $\lambda_f$ is \textbf{countercyclical}: tight markets $\to$ harder to fill vacancies.
\end{frame}

% ============================================================================
% SLIDE 18: SUMMARY
% ============================================================================
\begin{frame}{Summary and Takeaways}
    \begin{columns}[T]
        \begin{column}{0.48\textwidth}
            \textbf{Problem 1}:
            \begin{itemize}
                \item Chebyshev excels on smooth functions; splines handle kinks and local features
                \item 3D: cubic spline most accurate; Chebyshev competitive on smooth regions
                \item Occupational choice + collateral constraints drive wealth concentration
                \item Entrepreneurial returns generate Pareto-like wealth tails
            \end{itemize}
        \end{column}
        \begin{column}{0.48\textwidth}
            \textbf{Problem 2}:
            \begin{itemize}
                \item DMP: two-equation system in $(S, \lambda_w)$ captures search frictions
                \item Newton converges fastest ($\sim$6 its); iterative is simple but slow ($\sim$125)
                \item Linearization accurate near steady state; misses tail curvature
                \item Amplification mechanism: productivity $\to$ surplus $\to$ tightness $\to$ labor-market outcomes
            \end{itemize}
        \end{column}
    \end{columns}
    \vspace{0.5em}
    \centering
    \small
    Code: \texttt{notebooks/pset2\_problem1.ipynb} and \texttt{notebooks/pset2\_problem2.ipynb}.
\end{frame}

% ============================================================================
% APPENDIX
% ============================================================================
\appendix

\begin{frame}[plain]
    \begin{center}
        {\Large \textbf{Appendix}}
    \end{center}
\end{frame}

% ============================================================================
% APPENDIX A: CALIBRATION TABLES
% ============================================================================
\begin{frame}{Appendix: Problem 1 --- Calibration}
    \footnotesize
    \begin{center}
    \begin{tabular}{llcl}
        \toprule
        \textbf{Parameter} & \textbf{Symbol} & \textbf{Value} & \textbf{Description} \\
        \midrule
        \multicolumn{4}{l}{\emph{Static occupational choice (Problem 1.2)}} \\
        Capital share & $\alpha$ & $1/3$ & Cobb--Douglas \\
        Returns to scale & $\eta$ & $0.85$ & Span-of-control \\
        Wage / Rental rate & $W$\,/\,$R$ & $1.0$\,/\,$0.10$ & \\
        Leverage & $\lambda$ & $2.0$ & Collateral multiplier \\
        \midrule
        \multicolumn{4}{l}{\emph{Dynamic extension}} \\
        $\beta$, $\gamma$, $r$ & & $0.95$, $2.0$, $0.02$ & CRRA preferences \\
        $e$: $n_e$, $\rho_e$, $\sigma_e$ & & $5$, $0.90$, $0.20$ & Labor ability \\
        $z$: $n_z$, $\rho_z$, $\sigma_z$ & & $5$, $0.90$, $0.30$ & Entrep.\ ability \\
        Asset grid & $n_a$, $\bar a$ & $120$, $200$ & Grid \\
        \bottomrule
    \end{tabular}
    \end{center}
\end{frame}

% ============================================================================
% APPENDIX B: NEWTON JACOBIAN
% ============================================================================
\begin{frame}{Appendix: Problem 2 --- Newton Jacobian Structure}
    Stack unknowns $x = (S_1,\ldots,S_N,\lambda_{w,1},\ldots,\lambda_{w,N})'$. The $2N\times 2N$ Jacobian has block structure:
    \[
    J = \begin{pmatrix}
    \partial F^S/\partial S & \partial F^S/\partial\lambda_w \\
    \partial F^\lambda/\partial S & \partial F^\lambda/\partial\lambda_w
    \end{pmatrix}
    \]
    \vspace{0.2em}
    \small
    \begin{align*}
    \frac{\partial F^S}{\partial S} &= I - \beta\,\mathrm{diag}(1-\sigma-\gamma\lambda_w)\,P \\[4pt]
    \frac{\partial F^S}{\partial\lambda_w} &= \gamma\beta\,\mathrm{diag}(PS) \\[4pt]
    \frac{\partial F^\lambda}{\partial S} &= -\frac{1-\alpha}{\alpha}\,\mathrm{diag}\!\left(\frac{\lambda_w}{PS}\right)P \\[4pt]
    \frac{\partial F^\lambda}{\partial\lambda_w} &= I
    \end{align*}
    \normalsize
    Newton converges quadratically; Broyden approximates $J$ with rank-one updates.
\end{frame}

\end{document}
